\begin{abstract}
本稿では，検出と命名を分離した粒度可変物体アノテーションにおいて，現在の構成（CLIP-ViT-B/32とソフトマックス，COCOの80クラスを深さ0とするハイブリッド構成）に対し，3つの変更を加えた構成を比較する．
第一の変更は，視覚言語モデル（Vision-Language Model，VLM）をSigLIPに置き換え，候補語彙の中で正規化するソフトマックスの代わりに語ごとに独立なシグモイドで確信度を計算することである．
第二の変更は，YOLOの検出ボックスのクロップに対するVLMの最大確信度がYOLOの検出スコアを上回る場合にVLMの語を深さ0とすることである．
第三の変更は，YOLOを再学習せずにCOCO外の語へ対応するため，WordNetにおいてCOCOのカテゴリと同じ深さにある語彙を深さ0の候補とすることである．
COCO 2017 Validationの重複しない検証用・テスト用の分割とLVISの細かい粒度の正解を用いた評価により，以下を確認した．
シグモイドは，深さ0が誤った枝における最深段の確信度の識別性能（AUROC）をハイブリッド構成で0.418から0.617へ改善する一方，深さ0が正しい場合の深さ1の正誤の識別ではソフトマックスに劣る（0.762に対し0.644）．
スコアの比較による切替は，ソフトマックスではCOCOのTop-1正解率を0.969から0.936（CLIP）および0.858（SigLIP）へ低下させ，シグモイドでは切替がほとんど生じない．
WordNetの同じ深さの語彙（681語）を用いると，COCO外の物体の深さ0が正しい上位語となる割合は0から3.8\%に増えるが，COCOのTop-1正解率は0.945へ低下する．
\end{abstract}

\noindent
{\bf キーワード：}SigLIP，シグモイド，ソフトマックス，確信度，WordNet，オープンボキャブラリ物体アノテーション
